\section{讲座主线：能力上升与开放性下降的剪刀差}

\subsection{开场判断：Why Access Matters}
Percy Liang 在开场几分钟就给出整场讲座的核心命题：\textbf{Foundation Model 能力快速上升，但研究者可获得的访问权限在同步下降}。这种“能力曲线向上、开放曲线向下”的剪刀差，不只是商业策略变化，而是会直接重塑未来五年的研究议程。

\begin{importantbox}{Access Shapes Research}
讲者用一句话概括：\textit{``Access shapes research.''}  
当你只有 API 访问时，你最容易做 Prompt Engineering；当你拥有权重和训练代码时，你才有机会做机制解释、训练改造、安全审计与架构创新。
\end{importantbox}

\begin{figure}[H]
\centering
\includegraphics[width=0.93\textwidth]{frame-01-title.jpg}
\caption{讲座开场页：从 Agent 话题引向开放科学问题}
\end{figure}
\noindent\footnotesize 画面证据：字幕约在 00:00:09--00:00:51，讲者明确提出 “importance of open source” 及 “access shapes research”。\normalsize

\subsection{历史视角：NLP 的三次“访问红利”}
讲者回顾了 NLP/ML 历史中几次典型跃迁：
\begin{enumerate}
    \item 90 年代到 00 年代：互联网文本可访问性提升，催生统计 NLP；
    \item 10 年代：标注数据平台成熟，推动监督学习与基准化评测；
    \item 深度学习时代：GPU 与集群算力可获得性提升，Transformer 规模化成为现实。
\end{enumerate}

\begin{knowledgebox}{课程层面的隐含方法论}
这段历史回顾并不是“科普插曲”，而是为后文铺垫一个方法论：  
\textbf{讨论模型能力时，必须同时讨论数据访问、计算访问和系统访问。}
\end{knowledgebox}

\begin{warningbox}{常见误区：把开放问题简化为“是否开源权重”}
开放问题不是二元开关。权重是否可下载只是一个维度。若训练数据、训练流程、评测脚本、系统日志与部署策略不可见，研究可复现性仍然不足，科学结论依旧脆弱。
\end{warningbox}

\subsection{本章小结}
本章建立了整场讲座的总问题：我们不是在争论“开源好不好”，而是在讨论“不同访问层级会允许什么研究、阻断什么研究”，以及这种结构性差异如何影响 Agent 与 Foundation Model 的长期进展。


